A complete proof of \(\mathrm{C2}\) **cannot be obtained from only \(\mathrm{A2}\) and \(\mathrm{A4}\) as stated**. The assumptions are enough to prove a narrower two-strategy NPV dominance claim, but they do **not** by themselves justify the full conclusion “unique economically rational strategy at societal scale” or “irreversible,” and they do not bound any independent “benefits of reversion.”

Below is the rigorous diagnosis, followed by a corrected proof for the strongest statement that does follow under an additional no-independent-benefit convention.

---

## 1. Why the theorem is not provable as stated

To compare adoption and reversion, one must define the relative NPV

\[
\Delta(d,\rho)
=
\operatorname{NPV}_{\mathrm{rev}}(d,\rho)
-
\operatorname{NPV}_{\mathrm{cont}}(d,\rho).
\]

A natural continuous-time comparison would be

\[
\Delta(d,\rho)
=
B(d,\rho)
-
S(d)
+
\int_{0}^{\infty} e^{-\rho t}
\bigl(\Pi_0-\Pi_{\mathrm{AI}}(d)\bigr)\,dt,
\]

where \(B(d,\rho)\) denotes the present value of any independent economic benefits of reversion, such as avoided regulatory penalties, avoided AI operating costs, liability reductions, subsidies, or other non-productivity benefits.

Assumptions \(\mathrm{A2}\) and \(\mathrm{A4}\) imply only that

\[
\Pi_{\mathrm{AI}}(d)\ge (1+\gamma)\Pi_0
\]

and

\[
S(d)\ge c d^2.
\]

Therefore,

\[
\Delta(d,\rho)
\le
B(d,\rho)
-
c d^2
-
\frac{\gamma \Pi_0}{\rho}.
\]

The sign of \(\Delta(d,\rho)\) is not determined unless \(B(d,\rho)\) is bounded. No such bound appears in \(\mathrm{A2}\) or \(\mathrm{A4}\).

Indeed, take

\[
\Pi_{\mathrm{AI}}(d)=(1+\gamma)\Pi_0,
\qquad
S(d)=c d^2,
\qquad
B(d,\rho)=2c d^2.
\]

Then \(\mathrm{A2}\) and \(\mathrm{A4}\) both hold, but

\[
\Delta(d,\rho)
=
2c d^2
-
c d^2
-
\frac{\gamma \Pi_0}{\rho}
=
c d^2
-
\frac{\gamma \Pi_0}{\rho}.
\]

Hence

\[
\Delta(d,\rho)>0
\]

whenever

\[
d>\sqrt{\frac{\gamma \Pi_0}{c\rho}}.
\]

So in this admissible model, reversion has **positive** NPV for sufficiently large \(d\), contradicting the claimed lock-in conclusion. Thus \(\mathrm{C2}\) is not derivable from \(\mathrm{A2}\) and \(\mathrm{A4}\) alone.

---

## 2. What can be proved under a stricter two-strategy model

If the intended model is that reversion has **no independent benefits** beyond returning to baseline productivity \(\Pi_0\), then a rigorous proof is possible. But this is an additional modeling convention, not a consequence of \(\mathrm{A2}\) and \(\mathrm{A4}\).

### Corrected theorem

Assume, in addition to \(\mathrm{A2}\) and \(\mathrm{A4}\), that:

1. Productivity is measured in annual economic-value units.
2. The only two available strategies are:
   - continue AI use at integration depth \(d\);
   - revert completely to the pre-AI baseline.
3. Continuing AI has annual productivity \(\Pi_{\mathrm{AI}}(d)\).
4. Reversion incurs a one-time cost \(S(d)\) and thereafter yields annual productivity \(\Pi_0\).
5. Reversion has no other benefits or costs.

Then for every \(d\ge 1\) and every \(\rho>0\), reverting has strictly negative NPV relative to continuing AI. Hence, within this two-strategy model, continuing AI is the unique NPV-maximizing strategy.

### Proof

Define the NPV of continuing AI as

\[
V_{\mathrm{cont}}(d,\rho)
=
\int_0^\infty e^{-\rho t}\Pi_{\mathrm{AI}}(d)\,dt
=
\frac{\Pi_{\mathrm{AI}}(d)}{\rho}.
\]

Define the NPV of reverting to the baseline as

\[
V_{\mathrm{rev}}(d,\rho)
=
-S(d)
+
\int_0^\infty e^{-\rho t}\Pi_0\,dt
=
-S(d)+\frac{\Pi_0}{\rho}.
\]

Therefore the relative NPV of reversion is

\[
\Delta(d,\rho)
=
V_{\mathrm{rev}}(d,\rho)
-
V_{\mathrm{cont}}(d,\rho).
\]

Substituting the above expressions gives

\[
\Delta(d,\rho)
=
-S(d)
+
\frac{\Pi_0-\Pi_{\mathrm{AI}}(d)}{\rho}.
\]

By \(\mathrm{A2}\),

\[
\Pi_{\mathrm{AI}}(d)\ge (1+\gamma)\Pi_0.
\]

Hence

\[
\Pi_0-\Pi_{\mathrm{AI}}(d)
\le
\Pi_0-(1+\gamma)\Pi_0
=
-\gamma \Pi_0.
\]

Therefore,

\[
\frac{\Pi_0-\Pi_{\mathrm{AI}}(d)}{\rho}
\le
-\frac{\gamma \Pi_0}{\rho}.
\]

By \(\mathrm{A4}\),

\[
S(d)\ge c d^2.
\]

Thus

\[
-S(d)\le -c d^2.
\]

Combining these inequalities,

\[
\Delta(d,\rho)
\le
-cd^2
-
\frac{\gamma \Pi_0}{\rho}.
\]

Since \(c>0\), \(d\ge 1\), \(\gamma>0\), \(\rho>0\), and \(\Pi_0\ge 0\), we have

\[
-cd^2
-
\frac{\gamma \Pi_0}{\rho}
<0.
\]

Therefore,

\[
\Delta(d,\rho)<0.
\]

Equivalently,

\[
V_{\mathrm{rev}}(d,\rho)
<
V_{\mathrm{cont}}(d,\rho).
\]

Thus reverting to the pre-AI baseline has strictly negative NPV relative to continued AI use for every \(d\ge 1\) and every \(\rho>0\).

In this restricted model, one may take

\[
d^*(\gamma,c,\rho)=1.
\]

Indeed, no nontrivial threshold is needed: reversion is already strictly dominated for all integration depths covered by the assumptions.

\(\square\)

---

## 3. If reversion benefits are allowed, an extra assumption is needed

A more general valid theorem would require a bound on reversion benefits. For example, suppose the present value of all independent benefits of reversion satisfies

\[
B(d,\rho)\le \overline B(\rho)
\]

for all \(d\ge 1\). Then

\[
\Delta(d,\rho)
\le
\overline B(\rho)
-
c d^2
-
\frac{\gamma \Pi_0}{\rho}.
\]

Thus \(\Delta(d,\rho)<0\) whenever

\[
c d^2
>
\overline B(\rho)
-
\frac{\gamma \Pi_0}{\rho}.
\]

Equivalently, it suffices that

\[
d>
\sqrt{
\frac{
\bigl(\overline B(\rho)-\gamma\Pi_0/\rho\bigr)_+
}{c}
},
\]

where

\[
x_+ := \max\{x,0\}.
\]

So one possible threshold is

\[
d^*(\rho)
=
\max\left\{
1,
\sqrt{
\frac{
\bigl(\overline B(\rho)-\gamma\Pi_0/\rho\bigr)_+
}{c}
}
\right\}.
\]

But this threshold depends on \(\overline B(\rho)\) and \(\Pi_0\). Therefore a threshold depending only on \(\gamma\), \(c\), and \(\rho\) is not justified by the stated assumptions.

---

## Conclusion

The theorem \(\mathrm{C2}\) is too strong as written. From \(\mathrm{A2}\) and \(\mathrm{A4}\) alone, one can rigorously prove only the following narrower result:

\[
\text{If reversion has no independent benefits, then reversion is NPV-negative for all } d\ge 1.
\]

The broader claims of societal-scale uniqueness and irreversibility require additional assumptions about aggregation, available strategies, bounded reversion benefits, and the absence of offsetting economic gains from abandoning AI.